% --- replaces the split description in the experimental setup ---
\paragraph{Unit of evaluation.} The 377 tactile images are not 377 independent observations.
The GelSight frames are captured as continuous bursts while the sensor is held against one
egg, and the acquisition time of every frame survives in its filename. Clustering frames by
timestamp (a gap above 30\,s starts a new burst) recovers 19 contact episodes, 10 of whole
eggs and 9 of cracked eggs (Table~\ref{tab:episodes}); within an episode successive frames
are a median 1.43\,s apart and differ only slightly. The effective sample size of the
dataset is therefore 19 eggs, not 377 images.

The released split assigns images at random, which places near-duplicate frames of the same
physical egg on both sides of the partition: 18 of the 19 episodes appear in more than one
split, and every one of the 38 test images comes from an episode that also contributes a
training or validation frame. (The export contains no byte-identical or augmented copies;
the overlap is between distinct but adjacent frames of one continuous contact.) A score
computed on that split measures recall of eggs the detector has already seen, and because
the 38 test images are spread over 16 already-seen eggs it also carries no usable estimate
of variance.

\paragraph{Grouped cross-validation.} We therefore re-evaluate every detector with 5-fold
cross-validation grouped by contact episode and stratified by class. Each fold holds out
3--4 complete episodes (63--84 images, both classes always present) for testing and draws a
disjoint set of 3--4 episodes from the remaining pool for checkpoint selection, leaving
238--250 images for training (Table~\ref{tab:folds}). Every episode is held out exactly
once, so the 5 test sets partition all 377 images and each image is scored by exactly one
model that never saw its egg.

We report two aggregates, because they answer different questions. The mean and standard
deviation over folds describe how much the score moves when the held-out eggs change, which
is the quantity the released single split cannot provide. The pooled score instead
concatenates the per-fold detections and computes one average precision over all 377 images
under a single confidence ranking; it is the stricter and more directly comparable number,
since averaging five small per-fold APs rewards folds that happen to contain few, easy
instances. The gap between the two aggregates is itself a reading of how little data each
fold holds.

\paragraph{Arms.} So that the protocol change is not confounded with a change of training
recipe, we run a crossed design: both the released image-level split and the grouped folds,
each with both the released initialisation (random weights, 350 epochs) and COCO-pretrained
weights fine-tuned for 100 epochs. Everything else -- architecture family, image size, batch
size, augmentation, seed -- is held fixed, and all four arms are scored by the same
evaluation code on the test split. The effect of the protocol is then read within a single
initialisation, and vice versa.
